\chapter{Будущее: стоит ли переходить на гиперграфы}
\label{ch:hypergraphs}

Это приложение --- рассуждение о структурном направлении, а не готовая
рекомендация к внедрению. Вопрос: логично ли в будущем сделать память Рин
\emph{гиперграфом}, а не обычным графом.

\section{Проблема, которая это мотивирует}

Обычный граф связывает узлы \emph{попарно}. Но вывод часто \emph{n-арен}: он
рождается только при одновременном присутствии нескольких посылок, и ни одно
подмножество не достаточно. Канонический пример (глава~\ref{ch:fusion}):
\[
A=\text{``есть коты''},\quad B=\text{``живёт в Португалии''},\quad
C=\text{``ездит на пары''},
\]
\[
\begin{aligned}
X=f(A,B,C)
&=\text{``купи котам корм на выходных}\\
&\qquad \text{(в будни по пути с пар не успеть)''}.
\end{aligned}
\]
Каждая посылка ещё и подтягивает свой контекст (коты$\to$нужен корм;
Португалия$\to$магазины рано закрываются; универ$\to$поздний возврат;
см.\ главу~\ref{ch:fusion}), и только их совмещение даёт $X$. Ни $A\wedge B$,
ни $B\wedge C$, ни $A\wedge C$ не дают $X$. Попарные рёбра между посылками
могут вообще отсутствовать: ``Португалия'' и ``коты'' напрямую не связаны.

\section{Почему симплексы --- плохой ответ}

Первое, что приходит в голову, --- симплициальные комплексы: тройка
$\{A,B,C\}$ как 2-симплекс (треугольник). Но у симплициального комплекса
есть жёсткое требование \emph{downward closure}: если есть 2-симплекс
$\{A,B,C\}$, то обязаны быть и все его грани --- рёбра $A-B$, $A-C$, $B-C$.
А в нашем случае этих попарных связей \emph{нет и не должно быть}:
``Португалия'' и ``коты'' не связаны. Симплициальная память вынудила бы нас
проводить \emph{ложные} попарные рёбра ради формальной корректности
структуры. Это плохой обмен.

\section{Гиперребро --- прямой ответ}

\emph{Гиперребро} связывает \emph{множество} узлов сразу:
$e=\{A,B,C\}$, и оно \emph{не требует} существования попарных рёбер между
ними. Это ровно та семантика, что нужна n-арному выводу: ``эти три вместе
порождают это'', без обязательства, что они связаны попарно.

\begin{intuition}[Гиперграф против графа]
В графе ``дружба'' --- всегда пара. В гиперграфе можно сказать ``эти пятеро
были в одной поездке'' одним ребром, не утверждая, что каждый дружит с
каждым. Для памяти это разница между ``связь между двумя воспоминаниями'' и
``условие из нескольких воспоминаний, порождающее вывод''.
\end{intuition}

\section{Ключевое наблюдение: мы уже наполовину гиперграф}

Важно не пропустить: \textbf{текущая схема Рин уже кодирует гиперрёбра} ---
просто не называя их так. Когда fusion порождает узел $X$ с провенансом
$\{A,B,C\}$ (глава~\ref{ch:fusion}), это factor-graph-кодировка гиперребра
$\{A,B,C\}\to X$: сам вывод $X$ становится узлом, а его рёбра-обоснования
указывают на посылки. Попарных рёбер между $A,B,C$ при этом не требуется.
Формально это просто \textbf{двудольный граф} (узлы-посылки $\leftrightarrow$
узлы-факторы), и его \emph{полностью достаточно} для n-арного вывода ---
никакой высшей размерности не нужно.

Значит ``переход на гиперграфы'' --- это не ``переписать всё'', а более узкий
вопрос: \emph{делать ли гиперрёбра первоклассными объектами} (со своим
обходом и спектральной теорией), или достаточно узла-фактора.

\section{Первоклассные гиперрёбра: за и против}

\paragraph{За.} Появляются \emph{гиперграф-нативные} операции:
обход/spreading activation, уважающий \emph{совместность} срабатывания
(гиперребро ``зажигается'', когда активны \emph{все} его узлы, а не любой
один); гиперграфовые лапласианы и спектральная диагностика (аналоги
главы~\ref{ch:topology} для n-арной структуры). Это может точнее
моделировать ``условие из нескольких посылок''.

\paragraph{Против.} Алгоритмика гиперграфов менее зрелая и дороже графовой.
На практике гиперграф часто \emph{сводят} к графу --- clique-expansion
(заменить гиперребро кликой, что как раз возвращает ложные попарные рёбра!)
или star-expansion (ввести узел-фактор --- то есть ровно наша
factor-кодировка). Есть спектральная теория гиперграфов (нормализованный
гиперграфовый лапласиан; тензорные методы), но она сложнее и вычислительно
тяжелее, а динамический мутирующий гиперграф --- тем более.

\section{Рекомендация}

\begin{principle}[Гиперграф --- да, но постепенно]
Ближний горизонт: оставить \emph{factor-узлы} (вывод $X$ с
провенансом-множеством). Это уже даёт n-арную семантику без гиперграфовой
машинерии и без ложных попарных рёбер. Средний горизонт: сделать
гиперрёбра первоклассными \emph{только там}, где нужны гиперграф-нативные
операции --- совместное срабатывание при обходе и n-арная спектральная
диагностика. Симплициальные комплексы обходить стороной: их downward closure
противоречит разреженности реальных выводов.
\end{principle}

Заметим, что star-expansion (факторные узлы) и первоклассные гиперрёбра ---
это два взгляда на одно и то же; переход между ними инкрементален, а не
революционен. Полное же объединение узлов, рёбер и выводов в объекты одного
рода разной размерности --- это уже клеточные комплексы
(приложение~\ref{ch:cellcomplex}).
